\documentclass[hangout]{beamer}
%\documentclass[compress]{beamer}

\usepackage[utf8]{inputenc} 
\usepackage{listings}
\usepackage[T1]{fontenc}
\usepackage{lmodern}
\usepackage{graphicx}
\usepackage{color}
\usepackage[french]{babel}

\usetheme{Warsaw}

\title[ES lab]{Cours elasticsearch} 
\author{Johnny Mari\'ethoz} 
\institute{RERO}
\logo{\includegraphics[width=0.15\textwidth]{img/logo_rero}}
\lstset{language=Python}


\lstdefinestyle{custompython}{
  language=Python,
  basicstyle=\footnotesize\ttfamily,
  keywordstyle=\bfseries\color{green!40!black},
  commentstyle=\itshape\color{purple!40!black},
  identifierstyle=\color{blue},
  stringstyle=\color{orange},
}

\lstset{style=custompython}

\date[Genève 2014]{Cours HESGE} 

\begin{document}

\maketitle

\begin{frame}{Table des matières I}
	\tableofcontents[part=1]
	\tableofcontents[part=2]
\end{frame}
\begin{frame}{Table des matières II}
	\tableofcontents[part=3]
	\tableofcontents[part=4]
\end{frame}

\part{1}
\begin{frame}{Moteurs de recherche}
	\tableofcontents[part=1]
\end{frame}
\section{Moteurs de recherche: la partie visible}

\begin{frame}
	\begin{block}{Interface utilisateur}
	\begin{itemize}
	    \item \structure{Champ de recherche} entrée du texte de recherche
	    \item \structure{Facettes} fréquence d'un \alert{terme}
	    \item \structure{Filtres} \alert{réduit} la liste des résultats (click sur les facettes)
	    \item \structure{Tris} ordonnancement des résultats
	    	\begin{itemize}
	    		\item par titre
	    		\item par auteur
	    		\item par \alert{pertinence}
	    		\item par date de publication
	    	\end{itemize}
	    \item \structure{Scoring} pondération des champs pour le tris par pertinence
	    \item \structure{Pagination} nombre de résultats par page
	\end{itemize}
	
	\end{block}
\end{frame}


\part{2}
\section{Théorie: Information Retrieval}
\begin{frame}{Théorie: Information Retrieval}
	\tableofcontents[part=2]
\end{frame}

\frame{
	\begin{block}<1->{Définition}
\begin{quote}
	La recherche d'information (RI1) est le domaine qui étudie la manière de
	\alert{retrouver des informations} dans un \alert{corpus}. [Wikipédia]
\end{quote}
\end{block}
	}

\subsection{Deux étapes}
\frame{
	\begin{block}<1->{Indexation}
		\begin{itemize}
			\item création d'un \alert{index de descripteurs}
			\item lors de l'ajout d'un document
			\item peut \alert{prendre du temps} (machine)
		\end{itemize}
	\end{block}
	\begin{block}<2->{Recherche}
		\begin{itemize}
			\item comparaison d'une \alert{requête} avec un lot de \alert{documents}
			\item doit être \alert{rapide}
		\end{itemize}
	\end{block}
}

\subsection{Principes}
\frame{
	\begin{block}<1->{Comment ça marche?}
Produit interne (Inner product):
$$
	\bordermatrix{
	&\mathbf{q}&\cr
    \mathbf{term_1}& 1      \cr
    \mathbf{term_2}& 0      \cr
    \mathbf{\vdots}& \vdots \cr
    \mathbf{term_n}& 0      
	}
\cdot
	\bordermatrix{
	&\mathbf{d_1}&\mathbf{d_2}&\ldots &\mathbf{d_n}\cr
        & 1      &  0     & \ldots & 0 \cr
        & 0      &  1     & \ldots & 0 \cr
        & \vdots & \vdots & \ddots & \vdots\cr
        & 0      &  0     &\ldots  & 0
	}
$$

\structure{Note}: dans la vraie vie les vecteurs sont sparses et on ajoute des
poids: \url{http://en.wikipedia.org/wiki/Vector_space_model}
	\end{block}
}

\frame{
	%\frametitle{Aim}
	\begin{block}<1->{Principe de génération des descripteurs (termes)}
		\begin{itemize}
			\item pré-traitements
				\begin{itemize}
					\item \alert{stemming}: racinisation ou désuffixation
					\item \alert{tokenization}: décomposition en mots
				\end{itemize}
			\item extraction des descripteurs sous forme de vecteur d'entiers de la taille du nombre de descripteurs
			\item rajout de règles empiriques, recherche n-gram, synonymes, etc.
			\item un temps de réponse faible => optimisation + vecteur sparse (Lucene)
		\end{itemize}
	\end{block}
}

\subsection{Mesures}
\frame{
	%\frametitle{Aim}
	\begin{block}<1->{Mesures de performance}
		\begin{itemize}
			\item "precision and recall"
				\begin{itemize}
					\item on fixe un nombre de résultat (par ex. 10) et on \alert{compte} le nombre de résultats pertinents
				\end{itemize}
			\item demande des données \alert{annotées}
		\end{itemize}
	\end{block}
}


\subsection{Limitations}
\frame{
	%\frametitle{Aim}
	\begin{block}<1->{Limitations}
		\begin{itemize}
			\item des avions: stop word
			\item pain (boulangerie), pain(english)
			\item thé - the
			\item \alert{qualité} des données
		\end{itemize}
	\end{block}
}

\part{3}
\begin{frame}{Elasticsearch}
	\tableofcontents[part=3]
\end{frame}
\section{Elasticsearch}

\subsection{Concepts}
\frame{
	%\frametitle{Aim}
	\begin{block}{Concepts similaires à SQL}
		\begin{description}
			\item[index] base de donnée
			\item[document type] table
			\item[document] ligne d'une table
			\item[term] colonne
		\end{description}
	\end{block}
}

\subsection{"Analyzer" et "Mapping"}
\begin{frame}[fragile]
	%\frametitle{Aim}
	\begin{block}{Configuration de l'index (avant indexation)}
		\begin{lstlisting}
{
    "analyzer" : {
        "custom_analyzer" : {
            "tokenizer" : "keyword",
            "filter" : ["standard", "asciifolding", 
                        "lowercase"]
        }
    }
}
		\end{lstlisting}
	\end{block}
\end{frame}

\begin{frame}[fragile]
	%\frametitle{Aim}
	\begin{block}{Mapping (avant indexation)}
		Ajout de connaissances "a priori"
		\begin{lstlisting}
 "language" : {
     "type": "string",
     "index": "not_analyzed"
 }
		\end{lstlisting}
		\begin{itemize}
			\item type de données (integer, string, etc.)
			\item "index": "not\_analyzed" == "analyzer":"keyword"
		\end{itemize}
	\end{block}
\end{frame}

\begin{frame}[fragile]
	%\frametitle{Aim}
	\begin{block}{Mapping complèxe}
		\begin{lstlisting}
 "authors" : {
     "properties" : {
         "other_authors": {
             "type": "string",
             "fields" : {
                 "other_authors": {
                     "copy_to" : ["facet_authors"],
                     "type": "string",
                     "analyzer": "simple"
                 } } } } 
},
 "facet_authors" : {
     "type" : "string",
     "index": "not_analyzed"
 }
		\end{lstlisting}
	\end{block}
\end{frame}

\subsection{Query}
\begin{frame}[fragile]
	%\frametitle{Aim}
	\begin{block}<1->{Structure d'une requête}
			\begin{lstlisting}
{
    "size": 20,
    "from": 0,
    "query": {
        "filtered": {
            "query": {},
            "filter": {}
        } 
    },
    "aggs": {},
    "highlight": {},
    "sort": []
}
			\end{lstlisting}
	\end{block}
\end{frame}


\subsection{Facettes et filtres}
\begin{frame}[fragile]
	%\frametitle{Aim}
	\begin{block}{Facets / aggregators (lors de la recherche)}
		\begin{lstlisting}
"aggs": {
    "language": {
        "terms": {
            "field": "language",
            "size": 10,
            "order": {
                "_count": "desc"
            }
        }
    }
}
		\end{lstlisting}
		\begin{itemize}
			\item auteurs, type de documents, etc.
			\item approximation du nombre d'occurrences
		\end{itemize}
	\end{block}
\end{frame}

\frame{
	%\frametitle{Aim}
	\begin{block}<1->{Filtres (recherche)}
		\begin{itemize}
			\item comme une requête \alert{booléenne}
			\item plus \alert{rapide} que "query"
			\item utilisé par les facettes
		\end{itemize}
	\end{block}
}

\subsection{Query\_string}
\begin{frame}[fragile]
	%\frametitle{Aim}
	\begin{block}<1->{Requête utilisateur (texte de recherche)}
		\begin{itemize}
			\item Lucene query (support de title:)
			\item boosting (pondération des termes, par ex. un titre a plus de poids qu'une note)
		\begin{lstlisting}
"query_string": {
    "fields" : ["title.title^4", 
                "title.subtitle^2",
                "abstract"],
    "query": user_query
}
		\end{lstlisting}
		\end{itemize}
	\end{block}
\end{frame}

\subsection{Autres}
\frame{
	%\frametitle{Aim}
	\begin{block}<1->{"Highlights" (recherche)}
		\begin{itemize}
			\item \alert{mise en évidence} du descripteur correspondant à la recherche
			\item supporte le \alert{"stemming"}
		\end{itemize}
	\end{block}
}

\begin{frame}[fragile]
	%\frametitle{Aim}
	\begin{block}<1->{Tris (recherche)}
		\begin{itemize}
			\item tris par terme ou par \alert{pertinence} (valeur par défaut)
		\begin{lstlisting}
"sort" : [
    { 
	    "title.title" : {
            "order" : "asc"
        }
    }
]
		\end{lstlisting}
		\end{itemize}
	\end{block}
\end{frame}

\part{4}
\begin{frame}{En pratique}
	\tableofcontents[part=4]
\end{frame}
\section{En pratique}
\frame{
	%\frametitle{Aim}
	\begin{block}<1->{Outil: \alert{"es\_lab"}}
		\begin{itemize}
			\item utilisation de l'extension "sense"
			\item basé sur \alert{python}-Flask
			\item conversion des données MarcXML en \alert{JSON}
			\item script d'indexation
			\item configuration des indexes et des termes
				\begin{itemize}
					\item marc2json
					\item facettes
					\item highlight
					\item mapping
				\end{itemize}
			\item interface utilisateur (page web)
		\end{itemize}
	\end{block}
}

\end{document}
